\documentclass[11pt]{article}
\usepackage[utf8]{inputenc}
\usepackage{amsmath, amssymb, amsthm}
\usepackage[margin=1in]{geometry}

\theoremstyle{definition}
\newtheorem{definition}{Definition}
\newtheorem{theorem}{Theorem}
\newtheorem{lemma}{Lemma}
\newtheorem{proposition}{Proposition}

\theoremstyle{remark}
\newtheorem{remark}{Remark}

\title{Lecture 8: Parabolic Evolution Equations, Galerkin Approximation, and Nonlinear Methods}
\date{04.12.2025}
\author{Tien Anh Vu}

\begin{document}
\maketitle

This lecture develops the basic functional analytic framework for parabolic evolution equations, from the Gelfand triple and coercive operators to Galerkin approximation, a priori estimates, monotonicity methods, and compactness arguments for nonlinear problems.

\section*{1. Parabolic Evolution Equations}

\subsection*{1.1. Building the Foundation: Function Spaces and the Gelfand Triple}
Let \(V\) and \(H\) be separable, real Hilbert spaces such that
\[
V\subseteq H.
\]

A standard example is
\[
V=H_0^{1,2}(D),
\qquad
H=L^2(D).
\]
In this setup, \(V\) is densely and continuously embedded into \(H\), which we write as
\[
V\hookrightarrow H.
\]
Because the embedding is continuous, the norm of an element in \(H\) is bounded by its norm in \(V\):
\[
\|u\|_H\le \|u\|_V.
\]
Every element of \(V\) lies in \(H\), but the converse need not hold.

Identifying \(H\) with its dual \(H'\) via the Riesz representation theorem, we obtain the standard Gelfand triple
\[
V\hookrightarrow H\cong H'\hookrightarrow V'.
\]

\subsection*{1.2. Defining the Operator}
With the spaces fixed, we introduce an operator
\[
A:V\to V'.
\]
We require \(A\) to be a continuous, linear operator.

The key assumption is coercivity:
\[
-\langle Au,u\rangle_{V',V}\ge \alpha\|u\|_V^2-\lambda\|u\|_H^2
\]
for some constant \(\alpha>0\) and some real number \(\lambda\in\mathbb R\).

As a guiding example, consider the Laplacian operator
\[
Au=\Delta u.
\]
If we multiply this by a test function \(v\) and integrate over the domain \(D\), integration by parts yields
\[
-\int_D \Delta u\,v\,dx=\int_D \langle \nabla u,\nabla v\rangle\,dx.
\]
This shows how \(A\) acts through the duality pairing in \(V'\).

\subsection*{1.3. The Main Theorem: Existence and Regularity}
We can now state the basic well-posedness theorem.

\begin{theorem}
Let \(A\in L(V,V')\) be a coercive operator, let the initial data satisfy
\[
u_0\in H,
\]
and let the source term satisfy
\[
f\in L^2([0,T];V').
\]
Then the evolution equation
\[
\frac{du(t)}{dt}=Au(t)+f(t),
\qquad
t\in[0,T],
\]
with the initial condition
\[
u(0)=u_0,
\]
has a unique solution.

Moreover, this unique solution belongs to the intersection space
\[
u\in C([0,T];H)\cap L^2([0,T];V).
\]
\end{theorem}

This regularity statement has two parts:
\[
C([0,T];H)
\]
indicates that the solution is continuous in time, while
\[
L^2([0,T];V)
\]
indicates that the solution possesses space-time regularity measured in the stronger space \(V\).

For the heat equation
\[
\frac{du(t)}{dt}=\Delta u(t)+f(t),
\qquad
u(0)=u_0,
\]
the theorem guarantees a solution taking values in a first-order Sobolev space, namely
\[
u(t)\in H_0^{1,2}
\]
for almost every time \(t\).

\subsection*{1.4. Proving Uniqueness}
To prove uniqueness, let \(u\) and \(v\) be two solutions. Since the operator is linear, their difference \(u-v\) solves the homogeneous equation
\[
\frac{d}{dt}(u-v)=A(u(t)-v(t))\in V'.
\]
Since both solutions start from the same initial data \(u_0\), we also have
\[
(u-v)(0)=0.
\]

We track the \(H\)-norm of the difference and apply the identity
\[
\frac{d}{dt}\|u(t)-v(t)\|_H^2
=
2\left\langle \frac{d}{dt}(u-v)(t),(u-v)(t)\right\rangle_{V',V}.
\]
Substituting the differential equation into the right-hand side gives
\[
\frac{d}{dt}\|u(t)-v(t)\|_H^2
=
2\langle A(u-v)(t),(u-v)(t)\rangle_{V',V}.
\]
Coercivity makes the derivative non-positive. Since the initial difference is zero, we conclude that
\[
u(t)=v(t)
\qquad
\text{for all } t.
\]

\section*{2. Galerkin Approximation and Existence}

\subsection*{2.1. Concluding the Uniqueness Proof}
To finish the uniqueness argument, we use coercivity once more and obtain the bound
\[
-\alpha\|u-v\|_V^2,
\]
which is inherently less than or equal to zero.

Thus \(t\mapsto \|u(t)-v(t)\|_H\) is non-increasing. Since it starts at zero, it remains zero for all \(t\), and hence
\[
u(t)=v(t),
\]
proving uniqueness.

\subsection*{2.2. Proving Existence: The Galerkin Method}
For existence, we use the Galerkin method, replacing the infinite-dimensional problem by finite-dimensional approximations.

Choose a complete orthonormal system \(\{e_k\}\) in \(H\) with \(e_k\in V\), and define
\[
V_n:=\operatorname{span}\{e_1,\dots,e_n\}.
\]
Then
\[
V_n\subset V.
\]

\subsection*{2.3. Constructing the Finite-Dimensional ODE}
Projecting the original equation onto \(V_n\) yields a finite-dimensional linear ODE system. Let \(u_n(t)\) denote its unique solution, with
\[
u_n\in C([0,T];V_n).
\]

For each basis vector \(e_k\), the projected equation is
\[
\frac{d}{dt}\langle u_n(t),e_k\rangle_H
=
\langle Au_n(t),e_k\rangle_{V',V}
+
\langle f(t),e_k\rangle_{V',V}
\]
for \(t>0\). To ensure that this finite system mirrors the original problem, we also project the initial data:
\[
\langle u_n(0),e_k\rangle_H=\langle u_0,e_k\rangle_H.
\]

\subsection*{2.4. The Crucial Bridge: A Priori Energy Estimates}
To pass to the limit \(n\to\infty\), we need an a priori estimate, that is, a bound derived from the structure of the equation rather than an explicit formula for the solution.

The key estimate is the uniform bound
\[
\sup_{n\ge 1}
\left(
\sup_{t\in[0,T]}\|u_n(t)\|_H^2
+
\int_0^T \|u_n(t)\|_V^2\,dt
\right)
<\infty.
\]
This prevents the approximations from blowing up and provides the compactness needed for the limit passage.

\section*{3. The A Priori Estimate and Passage to the Limit}

\subsection*{3.1. Contextualizing with the Heat Equation}
For intuition, keep the heat equation in mind. We seek a uniform bound of the form
\[
\sup_{n\ge 1}\int_0^T
\left(
\|\nabla u_n(t)\|_{L^2(D)}^2+\|u_n(t)\|_{L^2(D)}^2
\right)\,dt
<\infty.
\]
This explains why Sobolev spaces are the natural choice for \(V\): weak derivatives are needed in order to control gradient terms.

\subsection*{3.2. Proving the A Priori Estimate: Differentiating the Norm}
Start with the time derivative of
\[
\frac12\|u_n(t)\|_H^2.
\]
Since \(u_n(t)\) lives in the finite-dimensional space spanned by the basis \(\{e_k\}\), we can write
\[
\frac12\frac{d}{dt}\|u_n(t)\|_H^2
=
\frac12\frac{d}{dt}\sum_{k=1}^n \langle u_n(t),e_k\rangle_H^2.
\]
Applying the chain rule to this finite sum yields
\[
\sum_{k=1}^n
\frac{d}{dt}\langle u_n(t),e_k\rangle_H\,
\langle u_n(t),e_k\rangle_H.
\]
Substituting the finite-dimensional evolution equation gives
\[
\sum_{k=1}^n
\langle Au_n(t)+f(t),e_k\rangle\,
\langle u_n(t),e_k\rangle_H.
\]
Because \(u_n(t)\) is a linear combination of these basis vectors, Parseval's identity collapses the sum into
\[
\langle Au_n(t)+f(t),u_n(t)\rangle.
\]
By linearity, this becomes
\[
\langle Au_n(t),u_n(t)\rangle+\langle f(t),u_n(t)\rangle.
\]

\subsection*{3.3. Applying Coercivity and Young's Inequality}
Coercivity gives
\[
\langle Au_n(t),u_n(t)\rangle\le -\alpha\|u_n(t)\|_V^2.
\]
For the forcing term, we apply the dual Cauchy-Schwarz inequality:
\[
\langle f(t),u_n(t)\rangle
\le
\|f(t)\|_{V'}\|u_n(t)\|_V.
\]
Combining both estimates gives
\[
\frac12\frac{d}{dt}\|u_n(t)\|_H^2
\le
-\alpha\|u_n(t)\|_V^2+\|f(t)\|_{V'}\|u_n(t)\|_V.
\]
To separate the mixed product, we use Young's inequality:
\[
\|f(t)\|_{V'}\|u_n(t)\|_V
\le
\frac{1}{2\alpha}\|f(t)\|_{V'}^2+\frac{\alpha}{2}\|u_n(t)\|_V^2.
\]
Substituting this back into the estimate yields
\[
\frac12\frac{d}{dt}\|u_n(t)\|_H^2
\le
-\frac{\alpha}{2}\|u_n(t)\|_V^2+\frac{1}{2\alpha}\|f(t)\|_{V'}^2.
\]

\subsection*{3.4. Integration and Establishing the Bound}
Integrating the differential inequality from \(0\) to \(t\), and moving the negative \(V\)-norm term to the left-hand side, gives
\[
\frac12\|u_n(t)\|_H^2
+
\frac{\alpha}{2}\int_0^t \|u_n(s)\|_V^2\,ds
\le
\frac12\|u_n(0)\|_H^2+\frac{1}{2\alpha}\int_0^t \|f(s)\|_{V'}^2\,ds.
\]
Taking the supremum over \(t\in[0,T]\), and using the bound on the projected initial data, yields a finite constant depending only on \(u_0\), \(f\), and \(T\). Hence \((u_n)\) is uniformly bounded in \(n\).

\subsection*{3.5. Passing to the Limit via Weak Convergence}
Uniform boundedness lets us extract a weakly convergent subsequence, still denoted by \(u_n\), such that
\[
u_n\rightharpoonup u
\qquad
\text{weakly in } L^2([0,T];V)
\]
for some limit function \(u\).

Since \(A\) is linear and continuous, we also have
\[
Au_n\rightharpoonup Au
\qquad
\text{weakly in } L^2([0,T];V').
\]
To verify this, we test against a function \(v(t)\) and write
\[
\int_0^T \langle Au_n(t),v(t)\rangle\,dt
=
\int_0^T \langle A^*v(t),u_n(t)\rangle\,dt
\to
\int_0^T \langle A^*v(t),u(t)\rangle\,dt
=
\int_0^T \langle Au(t),v(t)\rangle\,dt.
\]
Thus the weak limit satisfies the original infinite-dimensional equation, completing the linear existence proof.

\section*{4. Generalization to the Nonlinear Case}

\subsection*{4.1. The Nonlinear Roadblock}
We now turn to the nonlinear case. In the linear setting, the a priori bound produced a weakly convergent subsequence \(u_n\rightharpoonup u\), and linearity identified the weak limit of \(A(u_n)\). This is the step that fails in the nonlinear setting.

We still assume that our operator
\[
A:V\to V'
\]
is continuous and coercive, but we now assume it is not necessarily linear.

The a priori estimate remains valid, so we can still extract a weakly convergent subsequence
\[
u_n\rightharpoonup u
\qquad
\text{in } L^2([0,T];V).
\]
Furthermore, because the sequence of operators is bounded, we can also extract a subsequence such that
\[
A(u_n)\rightharpoonup g
\qquad
\text{in } L^2([0,T];V')
\]
for some limit function \(g\).

Weak convergence does not, in general, commute with nonlinear operators. Therefore, we cannot immediately conclude that
\[
A(u)=g.
\]

\subsection*{4.2. Strategy 1: Monotonicity (Dissipativity)}
The first way to bridge this gap is monotonicity, also referred to as dissipativity.

We impose three additional conditions on the nonlinear operator \(A\):

\begin{enumerate}
\item The monotonicity condition
\[
\langle A(u)-A(v),u-v\rangle_V\le \lambda\|u-v\|_H^2
\]
for some \(\lambda\in\mathbb R\).
\item Boundedness
\[
\|A(u)\|_{V'}\le M(1+\|u\|_V).
\]
\item A continuity assumption on the operator.
\end{enumerate}

By shifting the operator, we may absorb the \(\lambda\)-term and reduce the monotonicity inequality to one with right-hand side less than or equal to zero.

This gives the fundamental integral inequality
\[
\int_0^T \langle A(u_n)-A(v),u_n-v\rangle(t)\,dt\le 0
\]
for every test function \(v\in L^2([0,T];V)\).

\subsection*{4.3. Expanding and Passing to the Limit}
Expanding the pairing gives
\[
\langle A(u_n),u_n\rangle(t)-\langle A(u_n),v\rangle(t)-\langle A(v),u_n\rangle(t)+\langle A(v),v\rangle(t).
\]
Passing to the limit term by term, and using
\[
A(u_n)\rightharpoonup g,
\]
we obtain the convergence of \(\langle A(u_n),v\rangle\) to \(\langle g,v\rangle\). Likewise, because
\[
u_n\rightharpoonup u,
\]
the term \(\langle A(v),u_n\rangle\) converges to \(\langle A(v),u\rangle\). The final term is independent of \(n\). The delicate term is \(\langle A(u_n),u_n\rangle\), since both factors vary only weakly.

Because
\[
A(u_n)\rightharpoonup g,
\]
monotonicity supplies a limit-superior bound by \(\langle g,u\rangle\).

Substituting these limits back into the expanded integral, we arrive at
\[
\int_0^T \langle g-A(v),u-v\rangle(t)\,dt\le 0.
\]

\subsection*{4.4. Minty's Trick: Concluding \(A(u)=g\)}
To identify \(g\) with \(A(u)\), we use Minty's trick.

Choose the test function
\[
v=u+\varepsilon w,
\]
where \(w\) is arbitrary and \(\varepsilon>0\). Then
\[
u-v=-\varepsilon w,
\]
so the inequality becomes
\[
\int_0^T \langle g-A(u+\varepsilon w),w\rangle(t)\,dt\le 0.
\]

By the continuity assumption on \(A\), we may let \(\varepsilon\to 0\). This yields
\[
\int_0^T \langle g-A(u),w\rangle(t)\,dt\le 0.
\]
Because \(w\) is arbitrary, the only possible conclusion is
\[
g-A(u)=0.
\]
Therefore
\[
g=A(u).
\]

Thus the weak limit of the nonlinear operator applied to the approximations is exactly the operator applied to the limit solution, completing the existence proof under monotonicity.

\section*{5. The Browder-Minty Trick in Detail}

\subsection*{5.1. Managing the Nonlinear Inner Product}
We now examine the Browder-Minty argument more closely.

The main difficulty is passing to the limit in the term
\[
\int_0^T \langle A(u_n),u_n\rangle(t)\,dt.
\]
Because both factors in the pairing are weakly converging sequences, we cannot simply pass the limit inside.

To handle this, we return to the finite-dimensional Galerkin equation
\[
\frac{d}{dt}u_n(t)=A(u_n)(t)+f(t).
\]
Taking the inner product with \(u_n(t)\) and integrating from \(0\) to \(T\), we can rewrite the problematic term as
\[
\int_0^T \langle A(u_n),u_n\rangle(t)\,dt
=
\frac12\bigl(\|u_n(T)\|_H^2-\|u_n(0)\|_H^2\bigr)
-\int_0^T \langle f,u_n\rangle(t)\,dt.
\]
This reformulation replaces the nonlinear pairing by terms controlled through weak convergence and energy estimates.

\subsection*{5.2. Weak Lower Semicontinuity and the Limit Inferior}
Choose a further subsequence such that \(u_n(t)\rightharpoonup u(t)\) weakly, and pass to the limit in the rewritten identity.

The integral involving the fixed source term converges directly:
\[
\int_0^T \langle f,u_n\rangle(t)\,dt
\to
\int_0^T \langle f,u\rangle(t)\,dt.
\]
For the norm terms at the endpoints \(t=0\) and \(t=T\), we use the weak lower semicontinuity of the norm. This gives
\[
\frac12\bigl(\|u(T)\|_H^2-\|u(0)\|_H^2\bigr)
-\int_0^T \langle f,u\rangle(t)\,dt
\le
\liminf_{n\to\infty}
\left[
\frac12\bigl(\|u_n(T)\|_H^2-\|u_n(0)\|_H^2\bigr)
-\int_0^T \langle f,u_n\rangle(t)\,dt
\right].
\]
The left-hand side is the limit expression associated with \(g\), so we obtain
\[
\int_0^T \langle g(t),u(t)\rangle\,dt
\le
\liminf_{n\to\infty}
\int_0^T \langle A(u_n),u_n\rangle(t)\,dt.
\]
This is the key liminf estimate.

\subsection*{5.3. Assembling the Browder-Minty Inequality}
Returning to the monotonicity condition, for every test function \(v\) we know that
\[
\int_0^T \langle A(u_n)-A(v),u_n-v\rangle(t)\,dt\le 0.
\]
Expanding this pairing, passing to the weak limits \(A(u_n)\rightharpoonup g\) and \(u_n\rightharpoonup u\), and inserting the liminf bound from the previous step, we arrive at
\[
\int_0^T \langle g-A(v),u-v\rangle(t)\,dt
\le
\liminf_{n\to\infty}
\int_0^T \langle A(u_n)-A(v),u_n-v\rangle(t)\,dt
\le 0.
\]
This is the central inequality of the Browder-Minty argument.

\subsection*{5.4. Executing the Browder-Minty Trick}
Now choose
\[
v=u+\varepsilon w,
\]
where \(w\) is arbitrary and \(\varepsilon>0\). Then
\[
u-v=-\varepsilon w.
\]
Substituting this into the inequality and dividing by the positive constant \(\varepsilon\), we are led to the expression
\[
\int_0^T \langle g-A(u+\varepsilon w),w\rangle(t)\,dt\le 0.
\]
To pass to the limit \(\varepsilon\downarrow 0\), we use the continuity assumption on \(A\), which gives
\[
A(u+\varepsilon w)\to A(u)
\qquad
\text{in } V'.
\]
Hence the inequality simplifies to
\[
\int_0^T \langle g-A(u),w\rangle(t)\,dt\le 0.
\]

\subsection*{5.5. The Final Symmetry Argument}
Finally, the inequality
\[
\int_0^T \langle g-A(u),w\rangle(t)\,dt\le 0
\]
holds for every direction \(w\in L^2([0,T];V)\).

Since it holds for every \(w\), it also holds for \(-w\). Substituting \(-w\) yields
\[
-\int_0^T \langle g-A(u),w\rangle(t)\,dt\le 0,
\]
or equivalently
\[
\int_0^T \langle g-A(u),w\rangle(t)\,dt\ge 0.
\]
Hence the same quantity is both less than or equal to zero and greater than or equal to zero, so
\[
\int_0^T \langle g-A(u),w\rangle(t)\,dt=0
\qquad
\text{for all } w.
\]
By the fundamental lemma of the calculus of variations, this implies
\[
g=A(u)
\]
in the space \(L^2([0,T];V')\). This identifies \(g\) with \(A(u)\) and closes the nonlinear existence proof.

\section*{6. Strategy 2: Compactness}

\subsection*{6.1. The Limits of Monotonicity and the Need for Strategy 2}
We next ask what happens when the operator is not monotone.

The monotonicity argument is powerful, but it does not apply to some important equations from mathematical physics.

The main example is fluid dynamics, specifically the Navier-Stokes equations. Their nonlinear convective term destroys monotonicity, so the Browder-Minty trick cannot be applied directly. A different approach is needed.

This leads to Strategy 2: compactness.

\subsection*{6.2. The Power of Compact Embeddings}
The starting point is again the Gelfand triple
\[
V\hookrightarrow H\hookrightarrow V'.
\]
Here continuity is not enough; we need compactness.

Consider again the concrete example
\[
V=H_0^{1,2}(D),
\qquad
H=L^2(D).
\]
A foundational theorem in functional analysis, the Rellich-Kondrachov theorem, tells us that for bounded domains, the embedding
\[
H_0^{1,2}(D)\hookrightarrow L^2(D)
\]
is not only continuous but compact.

This matters because bounded sequences in compactly embedded spaces admit strongly convergent subsequences in the larger space. Strong convergence is the decisive tool for nonlinear problems.

\subsection*{6.3. Space-Time Compactness for the Galerkin Sequence}
Applying this to the Galerkin approximations \((u_n)\), the a priori estimate shows that they are bounded in the space-time space
\[
L^2([0,T];V).
\]
Because the spatial embedding \(V\hookrightarrow H\) is compact, this compactness can be transferred to the corresponding space-time setting under suitable hypotheses, for example through the Aubin-Lions lemma.

As a consequence,
\[
L^2([0,T];V)
\]
embeds compactly into
\[
L^2([0,T];H).
\]

\subsection*{6.4. Bypassing the Nonlinear Roadblock}
This changes the situation. Since \((u_n)\) is bounded in \(L^2([0,T];V)\), we can extract a subsequence that converges strongly to the limit \(u\) in
\[
L^2([0,T];H).
\]

This is what we need. If \(u_n\rightharpoonup u\) only weakly, then a continuous nonlinear operator \(A(u_n)\) need not converge to \(A(u)\). But if
\[
u_n\to u
\qquad
\text{strongly in } L^2([0,T];H),
\]
and the nonlinear operator satisfies suitable continuity assumptions with respect to the \(H\)-norm, then we may pass the limit directly through the operator and conclude that
\[
A(u_n)\to A(u).
\]
Thus compactness upgrades weak convergence to strong convergence and provides a second route for nonlinear evolution equations, including models where monotonicity fails.

\section*{Appendix}

\subsection*{A.1. Gelfand Triple}
\begin{definition}
A Gelfand triple is a chain of spaces
\[
V\hookrightarrow H\cong H'\hookrightarrow V',
\]
where \(V\) is densely and continuously embedded into the Hilbert space \(H\), and \(V'\) is the dual of \(V\).
\end{definition}

\subsection*{A.2. Coercive Operator}
\begin{definition}
An operator \(A:V\to V'\) is called coercive if there exist constants \(\alpha>0\) and \(\lambda\in\mathbb R\) such that
\[
-\langle Au,u\rangle_{V',V}\ge \alpha\|u\|_V^2-\lambda\|u\|_H^2
\qquad
\text{for all } u\in V.
\]
\end{definition}

\subsection*{A.3. Weak Form of the Laplacian}
\begin{definition}
The Laplacian can be defined weakly by the relation
\[
\langle \Delta u,v\rangle_{V',V}
:=
-\int_D \langle \nabla u,\nabla v\rangle\,dx,
\]
which makes sense whenever \(u,v\in V\).
\end{definition}

\subsection*{A.4. Evolution Equation}
\begin{definition}
An abstract parabolic evolution equation has the form
\[
\frac{du(t)}{dt}=Au(t)+f(t),
\qquad
u(0)=u_0,
\]
where \(A:V\to V'\), \(f\) is a forcing term, and \(u_0\) is the initial condition.
\end{definition}

\subsection*{A.5. Uniqueness by Energy Estimate}
\begin{lemma}
If the difference of two solutions satisfies a differential inequality forcing the \(H\)-norm of the difference to remain zero when it starts at zero, then the two solutions coincide.
\end{lemma}

\subsection*{A.6. Monotone Operator}
\begin{definition}
An operator \(A:V\to V'\) is called monotone if
\[
\langle A(u)-A(v),u-v\rangle_{V',V}\le \lambda\|u-v\|_H^2
\]
for all \(u,v\in V\) and some \(\lambda\in\mathbb R\). After shifting by \(\lambda I\), one often reduces to the case \(\lambda=0\).
\end{definition}

\subsection*{A.7. Galerkin Approximation}
\begin{definition}
Given a complete orthonormal system \((e_k)\) in \(H\) with \(e_k\in V\), the Galerkin subspaces are
\[
V_n=\operatorname{span}\{e_1,\dots,e_n\}.
\]
The Galerkin method replaces the infinite-dimensional evolution equation by its projection onto \(V_n\).
\end{definition}

\subsection*{A.8. A Priori Estimate}
\begin{definition}
An a priori estimate is a bound on approximate or exact solutions derived directly from the structure of the equation, independent of any explicit formula for the solution.
\end{definition}

\subsection*{A.9. Weak Lower Semicontinuity}
\begin{lemma}
If \(x_n\rightharpoonup x\) weakly in a Hilbert or Banach space, then
\[
\|x\|\le \liminf_{n\to\infty}\|x_n\|.
\]
\end{lemma}

\subsection*{A.10. Browder-Minty Trick}
\begin{remark}
The Browder-Minty trick is the standard monotonicity argument used to identify a weak limit \(g\) of \(A(u_n)\) with \(A(u)\). It combines monotonicity, a liminf estimate, and perturbations of the form \(v=u+\varepsilon w\).
\end{remark}

\subsection*{A.11. Compact Embedding}
\begin{definition}
An embedding \(V\hookrightarrow H\) is compact if every bounded sequence in \(V\) has a subsequence that converges strongly in \(H\).
\end{definition}

\subsection*{A.12. Rellich-Kondrachov Theorem}
\begin{theorem}
For bounded domains \(D\), the Sobolev embedding
\[
H_0^{1,2}(D)\hookrightarrow L^2(D)
\]
is compact.
\end{theorem}

\subsection*{A.13. Aubin-Lions Compactness Principle}
\begin{theorem}
Under suitable assumptions on the spaces and time derivatives, bounded subsets of \(L^2([0,T];V)\) are relatively compact in \(L^2([0,T];H)\) whenever the embedding \(V\hookrightarrow H\) is compact.
\end{theorem}

\end{document}
